%%% -*- coding: utf-8 -*-
% !TeX root = main-slides.tex

\lecture{SEE Concepts}{see-concepts}
\section{SEE Concepts}

\showtitlepage{Cities, Providers, Players, Actions, History}

%%% Every example is taken from SEE, so the sources named in the listings can
%%% be opened next to the slide.

\subsection{Code Cities}

\begin{frame}[fragile]{A Code City}
    \framelabel{slide:city}
    \begin{lstlisting}[basicstyle=\scriptsize\ttfamily]
// Assets/SEE/Game/City/AbstractSEECity.cs
public abstract partial class AbstractSEECity : SerializedMonoBehaviour
{
    public abstract Graph LoadedGraph { get; protected set; }
    public abstract IGraphRenderer Renderer { get; }
    public NodeTypeVisualsMap NodeTypes = new();
    ...
 }
    \end{lstlisting}
    \vspace{1ex}

    \begin{itemize}
        \item A \textit{code city} is a \texttt{GameObject} holding an
        instance of a subclass of \texttt{AbstractSEECity}.
        \item An \texttt{AbstractSEECity} has the settings to generate
        all \texttt{GameObject}s for a code city (loading and rendering a \texttt{Graph}).
        \item A \textit{game node} is a \texttt{GameObject} representing
        a graph \texttt{Node} in a code city.
        \item A \textit{game edge} is a \texttt{GameObject} representing
        a graph \texttt{Edge} in a code city.
    \end{itemize}
\end{frame}

\begin{frame}{\texttt{GameObject} Hierarchy of a Code City}
  \begin{columns}[c]
    \begin{column}{0.46\textwidth}
      \begin{itemize}
        \item There is at most one root game node that is an immediate child
        of a code city.
        \item The game-node hierarchy and graph-node hierarchy
        are isomorphic.
        \item All game edges are immediate children of the code city.
      \end{itemize}
    \end{column}
    \begin{column}{0.5\textwidth}
      %%% \incgraphics adds a \colorbox, whose \fboxsep would otherwise
      %%% push the image past the column edge.
      \incgraphics{width=\dimexpr\textwidth-2\fboxsep\relax}%
                  {figures/CloneCityHierarchy.png}
    \end{column}
  \end{columns}
\end{frame}

\begin{frame}[fragile]{Scene Constraints}
    \begin{itemize}
        \item  A scene may hold several cities side by side,
        each with its own data and settings.
        \item The \texttt{name} of a game node/edge equals the \texttt{ID}
        of a graph \texttt{Node}/\texttt{Edge}.
        \item The \texttt{name} of a game node/edge is unique within
        a scene at a particular point in time.
    \end{itemize}
    \vspace{1ex}

    \begin{lstlisting}[basicstyle=\scriptsize\ttfamily]
// game object -> graph element: read the NodeRef/EdgeRef component
gameObject.TryGetNode(out Node node);    // false if there is none
gameObject.TryGetEdge(out Edge edge);
Node n = gameObject.GetNode();           // throws if there is none

// graph element -> game object: by ID, which is the game object's name
GameObject gameNode = GraphElementIDMap.Find(node.ID);
GameObject gameEdge = GraphElementIDMap.Find(edge.ID);
    \end{lstlisting}
    \vspace{1ex}

    \hint{The way back leads through the \texttt{ID} only, and that is no
    accident: a net action (\slideref{slide:netaction}) cannot send a
    reference, but it can send an \texttt{ID}.}
\end{frame}


\begin{frame}[fragile]{The Family of Cities}
  \begin{lstlisting}[basicstyle=\scriptsize\ttfamily]
AbstractSEECity                  the component, the settings, the node types
 |
 +-- SEECity                     one graph
 |    +-- SEECityRandom          a generated graph, for experiments
 |    +-- SEEReflexionCity       implementation and architecture in one
 |    +-- VCSCity                anchored in a version control repository
 |         +-- BranchCity        the branches of a repository
 |         +-- CommitCity        the state at one commit
 |
 +-- SEECityEvolution            a series of graphs, played as an animation
  \end{lstlisting}
  \vspace{1ex}

  The split at the top is the one that matters: \texttt{SEECity} shows a
  single graph, \texttt{SEECityEvolution} a whole series of them. Everything
  below \texttt{SEECity} differs only in where the graph comes from and what
  extra structure it carries.
\end{frame}

\begin{frame}[fragile]{Settings Are Annotated, Not Coded}
  \begin{lstlisting}[basicstyle=\scriptsize\ttfamily]
// Assets/SEE/Game/City/SEECity.cs
[OdinSerialize, ShowInInspector,
 Tooltip("A graph provider yielding the data ..."),
 TabGroup(DataFoldoutGroup), RuntimeTab(DataFoldoutGroup), ...]
public SingleGraphPipelineProvider DataProvider = new();
  \end{lstlisting}
  \vspace{1ex}

  \begin{itemize}
    \item \texttt{TabGroup} places the field in the editor inspector,
      \texttt{RuntimeTab} in the configuration menu that SEE builds at run
      time by reflecting over these very fields.
    \item The same fields are written to and read from the city's
      configuration file by \texttt{ConfigWriter} and \texttt{ConfigReader}.
    \item One annotated field therefore yields three things at once: an
      inspector row, a menu entry, and a line in the configuration file. Adding
      a setting means adding a field, not writing user interface code.
  \end{itemize}
\end{frame}

\subsection{Graph Providers}

\begin{frame}[fragile]{Where the Graph Comes From}
  \framelabel{slide:provider}
  \begin{lstlisting}[basicstyle=\scriptsize\ttfamily]
// Assets/SEE/Game/City/SEECity.cs
public SingleGraphPipelineProvider DataProvider = new();

public virtual async UniTask LoadDataAsync()
{
  ...
  LoadedGraph = await DataProvider.ProvideAsync(new Graph(""), this,
                                                ReportProgress, token);
  ...
}
  \end{lstlisting}
  \vspace{1ex}

  The city itself reads no file and calls no tool. It hands an empty graph to
  its provider and takes back a finished one. Which files, which repository,
  which analysis --- none of that is the city's business.
\end{frame}

\begin{frame}[fragile]{A Graph Provider Transforms a Graph}
  \begin{lstlisting}[basicstyle=\scriptsize\ttfamily]
public abstract UniTask<T> ProvideAsync(T graph, AbstractSEECity city,
                                        Action<float> changePercentage,
                                        CancellationToken token);

// GXLSingleGraphProvider: ignores the input, produces a graph
return await GraphReader.LoadAsync(Path, city.HierarchicalEdges, ...);

// CSVGraphProvider: takes the graph, adds metrics, returns it
await MetricImporter.LoadCsvAsync(graph, Path, token: token);
return graph;
  \end{lstlisting}
  \vspace{1ex}

  Graph in, graph out --- that one signature is what makes providers
  composable. A provider is not a loader but a step, and only the first step
  of a pipeline usually ignores what it was given.
\end{frame}

\begin{frame}{Sources, Augmenters, Combiners}
  \begin{center}
    \footnotesize
    \begin{tabular}{ll}
      \hline
      \textbf{Role} & \textbf{For example}\\
      \hline
      Produce a graph      & \texttt{GXLSingleGraphProvider},
                             \texttt{GitBranchesGraphProvider}\\
      Add to a graph       & \texttt{CSVGraphProvider},
                             \texttt{JaCoCoGraphProvider},\\
                           & \texttt{ReportGraphProvider},
                             \texttt{DashboardGraphProvider}\\
      Combine graphs       & \texttt{ReflexionGraphProvider},
                             \texttt{MergeDiffGraphProvider}\\
      Hold other providers & \texttt{SingleGraphPipelineProvider}\\
      \hline
    \end{tabular}
  \end{center}
  \vspace{2ex}

  A typical pipeline reads a GXL file, then adds coverage from JaCoCo, then
  adds issues from a report. Each step is configured separately and none of
  them knows about the others.
  \vspace{1ex}

  \hint{The augmenters expect a graph and complain if they get none. Order in
  the pipeline is therefore not a matter of taste.}
\end{frame}

\begin{frame}[fragile]{One Graph or a Series}
  \begin{lstlisting}[basicstyle=\scriptsize\ttfamily]
public abstract class GraphProvider<T, K> where K : Enum

public abstract class SingleGraphProvider
  : GraphProvider<Graph, SingleGraphProviderKind>

public abstract class MultiGraphProvider
  : GraphProvider<List<Graph>, MultiGraphProviderKind>

// SEECity            LoadedGraph       = await ...(new Graph(""), ...)
// SEECityEvolution   LoadedGraphSeries = await ...(new List<Graph>(), ...)
  \end{lstlisting}
  \vspace{1ex}

  Evolution needs a sequence rather than a graph, and that is the only
  difference: the same framework, the same pipeline idea, a different
  \texttt{T}. The kind enum accompanying it lets a saved configuration name
  the provider to recreate.
\end{frame}

\subsection{Multiple Players}

\begin{frame}[fragile]{Every Change Happens Twice}
  \framelabel{slide:netaction}
  \begin{lstlisting}[basicstyle=\scriptsize\ttfamily]
// Assets/SEE/Controls/Actions/DeleteAction.cs
// Actual deletion (on all clients and locally).
new DeleteNetAction(go.name, removeNodeTypes).Execute();
GameElementDeleter.Delete(go, removeNodeTypes);

// Assets/SEE/Net/Actions/GraphElement/DeleteNetAction.cs
public override void ExecuteOnClient()
{
  GameObject objToDelete = Find(GraphElementID);
  GameElementDeleter.Delete(objToDelete, RemoveNodeTypes);
}
  \end{lstlisting}
  \vspace{1ex}

  The change itself lives in neither class. Both the action and its network
  twin call the same helper in \texttt{SEE.Game.SceneManipulation}, so what
  the other players see cannot drift from what you see. The local half is a
  full-blown action with a life cycle (\slideref{slide:action}); the remote
  half is one method.
\end{frame}

\begin{frame}[fragile]{The Life of a NetAction}
  \begin{lstlisting}[basicstyle=\scriptsize\ttfamily]
// Assets/SEE/Net/Actions/AbstractNetAction.cs
public void Execute(ulong[] recipients = null)
{
  Network.BroadcastAction(ActionSerializer.Serialize(this), recipients);
}

public virtual void ExecuteOnServer() { }   // default: nothing
public abstract void ExecuteOnClient();     // every client but the sender
  \end{lstlisting}
  \vspace{1ex}

  \begin{itemize}
    \item The action travels to the server as a string, and the server passes
      it on with \texttt{RpcTarget.Not(senderId)}: the sender has already made
      the change locally and must not make it twice.
    \item \texttt{Requester} carries the client identifier, so a receiver can
      tell whose action it is executing.
    \item Actions too large for one packet are cut into fragments and put back
      together on the far side.
  \end{itemize}
\end{frame}

\begin{frame}[fragile]{What May Travel}
  \begin{lstlisting}[basicstyle=\scriptsize\ttfamily]
// Assets/SEE/Net/Actions/GraphElement/DeleteNetAction.cs
// Note: All attributes are made public so that they will be serialized
// for the network transfer.
public bool RemoveNodeTypes;

// and never a reference; the object is found again on each client:
GameObject objToDelete = Find(GraphElementID);
  \end{lstlisting}
  \vspace{1ex}

  \begin{itemize}
    \item A net action is serialized with \texttt{JsonUtility}: every field
      that must travel is public and not \texttt{readonly}.
    \item A \texttt{GameObject} or a \texttt{Component} cannot travel, because
      the receiving client has different objects at different addresses. Send
      the identifier and look it up again.
    \item In the editor, \texttt{Execute} checks serializability before
      sending, so this mistake is caught while you are testing, not in a
      session with others.
  \end{itemize}
\end{frame}

\begin{frame}[fragile]{Late Joiners and Ephemera}
  \begin{lstlisting}[basicstyle=\scriptsize\ttfamily]
// Assets/SEE/Net/Actions/AbstractNetAction.cs
public virtual bool ShouldBeSentToNewClient { get => true; }

// Assets/SEE/Net/Actions/SetHoverNetAction.cs and its kin
public override bool ShouldBeSentToNewClient { get => false; }
  \end{lstlisting}
  \vspace{1ex}

  \begin{itemize}
    \item The server keeps every action that describes the world, and replays
      them to a client that joins later. That is how a newcomer arrives in a
      city somebody else built.
    \item Hovering, grabbing, highlighting, pointing, sounds: these describe a
      moment, not the world. They opt out, or a late joiner would inherit a
      history of somebody else's mouse.
    \item Deciding this correctly is part of writing a net action. Ask whether
      a player who arrives in an hour still needs to know.
  \end{itemize}
\end{frame}

\subsection{The Action History}

\begin{frame}[fragile]{An Action Has a Life Cycle}
  \framelabel{slide:action}
  \begin{lstlisting}[basicstyle=\scriptsize\ttfamily]
// Assets/SEE/Utils/History/IReversibleAction.cs
void Awake();            // once, when the action is first started
void Start();            // whenever it is (re-)enabled
bool Update();           // every frame; true means completed
void Stop();             // another action takes over
void Undo();  void Redo();
IReversibleAction NewInstance();

enum Progress { NoEffect, InProgress, Completed }
  \end{lstlisting}
  \vspace{1ex}

  An action is a mode, not a method call: it is started, then asked every
  frame whether it is done. When \texttt{Update} finally returns true, the
  history records the action and immediately puts a fresh instance of the same
  kind in its place, so the tool stays selected.
\end{frame}

\begin{frame}[fragile]{Two Histories}
  \begin{lstlisting}[basicstyle=\scriptsize\ttfamily]
// Assets/SEE/Utils/History/ActionHistory.cs
private Stack<IReversibleAction> UndoHistory { get; } = new();
private Stack<IReversibleAction> RedoHistory { get; } = new();

private readonly List<GlobalHistoryEntry> globalHistory = new();

public readonly struct GlobalHistoryEntry
{
  public bool IsOwner { get; }              // did I execute it?
  public HistoryType ActionType { get; }    // Action or UndoneAction
  public string ActionID { get; }
  public HashSet<string> ChangedObjects { get; }
}
  \end{lstlisting}
  \vspace{1ex}

  The two stacks are yours alone and hold the real objects. The global history
  is shared: every player has the same list, kept in step by
  \texttt{NetActionHistory} with \texttt{Push}, \texttt{Delete}, and
  \texttt{Replace}.
\end{frame}

\begin{frame}{Undoing in Company}
  \begin{itemize}
    \item Undo is not the last thing that happened but the last thing
      \emph{I} did. \texttt{IsOwner} in the global history is what makes that
      distinction possible, and \slideref{slide:netaction} is why the question
      arises at all.
    \item Before undoing, SEE looks at every entry recorded after yours. If
      somebody else has touched one of your \texttt{ChangedObjects} since,
      undoing would silently discard their work, and \texttt{UndoImpossible}
      is raised instead.
    \item An action that has had no effect yet is not undone but simply
      dropped: undoing it would give the user no visible answer to a key
      press.
  \end{itemize}
  \vspace{2ex}

  \hint{This is why \texttt{GetChangedObjects} is not optional paperwork. An
  action that under-reports what it changed makes conflicting undos possible;
  one that over-reports makes legitimate undos fail.}
\end{frame}

%%% Local Variables:
%%% mode: latex
%%% TeX-master: "main-slides"
%%% mode: reftex
%%% mode: flyspell
%%% End:
